\begin{center}
\textbf{요 약}
\end{center}

본 논문에서는 격자 기반 전자서명 \code{HAETAE}의 전체 서명 경로를
상류 연산, 응답 연산 및 거부 판정으로 구분하고, 소프트웨어 오류
에뮬레이션을 이용하여 각 지점의 취약 조건과 보안 영향을 분석한다.
특히 응답에서 비밀키 곱 항이 제거되는 T1과 일회용 난수 항이 제거되는
T2의 대수적 복원 관계를 도출하였다. T2 소프트웨어 실험에서는 30회 중
25회 오류 응답이 방출되었으며, 방출된 모든 경우에 서명 비밀 성분
$\svec_1$의 768개 계수를 복원하였다. T1에 대해서는 곱 결과 버퍼를
사전 초기화한 인과대조 구현에 실제 클럭 글리치를 주입하였다. 공격
구간을 식별한 후 정상 기준 응답과 두 개의 유용한 결함 응답을 이용하여
$\svec_1$ 전 계수를 복원하였고, 별도 실행에서도 594회 이내에 동일한
결과를 재현하였다. 한편 미수정 융합 응답 구현의 T2 탐색에서는 총
650회 동안 깨끗한 T2가 관측되지 않았다. 이는 소프트웨어로 정의한
오류 효과의 취약성과 실제 물리적 실현 가능성을 구분해야 하며, 버퍼
상태와 컴파일된 반복 구조가 공격 결과에 영향을 줄 수 있음을 보인다.
대응기법은 분석에서 도출한 최소 보안 요구사항과 후속 연구 방향으로
제시한다.

\vspace{0.4em}
\noindent\textbf{Keywords:} HAETAE, Lattice-based Signature, Fault Injection,
Instruction Skip, Clock Glitch, Secret Recovery
